\section*{Limitations}
The evaluation is retrospective and curated: the 100 retrieval questions and 50 paired generation cases do not form a paper-disjoint, independently sampled test set. The 150-case LLM-evaluation sample is disjoint from a reserved 50-case annotation packet, but both come from the same legacy 200-question, 10-paper corpus. Its source result file lacks the generator digest recorded for the paired study; its retained judge label does not establish the named model in the planning file; and one eligible case has a contradictory not-judgeable reason. We exclude its visual-grounding counts because the completed rows use a category absent from the frozen rubric. No LLM label is an independent human judgment. At the September 21, 2026 freeze, the human study has 250 planned but zero submitted ratings, and is itself development-overlapping rather than paper-disjoint. Automatic lexical overlap can reward copying, miss paraphrases, and overlook wrong numbers or unsupported implications; a failed lexical negative control prevents us from treating verifier labels as semantic support. Cross-paper phrase absence is not proof that a question is semantically unanswerable. Pipeline status is not equivalent to a refusal in prose. Page alignment is observable for only 20 of 30 answerable cases per condition because visual gold loci name figures or tables without page labels. The results use one pinned local model and one 18\,GB workstation; memory observations are pre-window free-memory checks rather than peak process measurements, and latency includes uncontrolled cache effects. The prototype has not been tested in a real researcher workflow, and no user benefit, production robustness, throughput, or cross-domain generalization is established. PDF extraction and reading order remain potential error sources. Local analysis does not imply that initial model or document acquisition is offline. The system should not be used for autonomous high-stakes decisions; users must inspect primary evidence.
